\begin{center}
    {\large \textbf{Giorno 4} \textbar\ Tagliatele la testa! \textbar\ 12 Agosto 2024 \textbar\ \textbf{Gili Asahan}, Lombok}
\end{center}

\noindent\rule{\textwidth}{0.4pt}
\begin{figure}[h!] % Portrait images below
    \centering
    % First portrait image (on the left)
    \begin{minipage}[h]{0.48\textwidth} % Set width equal to half the page
        \centering
        \includegraphics[width=\textwidth]{images/flags/Screenshot Gili A.png}
    \end{minipage}
    \hfill
    % Text
    \begin{minipage}[h]{0.48\textwidth}
        \raggedright
        \textbf{Superficie}: 4.567 km² \\
        \vspace{0.3cm}
        \textbf{Popolazione}: 3,3 milioni \\
        \vspace{0.3cm}
        \textit{L'isola}
    \end{minipage}
\end{figure}

\landscapeLargeMargin{images/day4/PXL_20240812_004417818.jpg}

\landscapeLargeMargin{images/day4/PXL_20240812_024559153.jpg}

\landscapeLargeMargin{images/day4/PXL_20240812_042505497.jpg}

\landscapeLargeMargin{images/day4/PXL_20240812_042543818.jpg}

\landscapeLargeMargin{images/day4/PXL_20240812_044451696.jpg}

\landscapeLargeMargin{images/day4/PXL_20240812_051220562.jpg}

\landscapeLargeMargin{images/day4/PXL_20240812_095840110.MP.jpg}

\noindent 12 Agosto 2024 \newline

Come previsto, alle quattro del mattino siamo svegliati dal Muezzin. Questo è l'orario giusto per le riflessioni filosofiche, e questa volta l'oggetto delle mie elucubrazioni sono i mass shooter — gli omicidi di massa, fenomeno purtroppo molto frequente negli Stati Uniti — e in particolare il recente trend di pazzi che entrano nelle moschee armati e aprono il fuoco su quante più persone riescano. Chiaramente non giustifico, ma se queste sparatorie sono avvenute alle quattro del mattino sicuramente rimuoverei la premeditazione dalle accuse e valuterei tutta una serie di attenuanti. In questo momento, dopo che per la seconda notte sono stato svegliato da preghiere arabe, anche io non mi sento più tanto sano di mente. Trovandomi in un villaggio di nativi a Lombok il massimo che potrei fare sarebbe lanciare qualche gallina contro la moschea, ma dubito servirebbe a qualcosa.

Per fortuna al termine delle preghiere riesco a riaddormentarmi e tiro fino all'ora di colazione. A Elisa non va altrettanto bene: anche lei è rimasta sveglia tutta la notte e quando si sveglia sembra la Regina di Cuori di Alice nel Paese delle Meraviglie, solo che ce l'ha con le galline. "Tagliatele la testa!" Inoltre, per non sentire il Muezzin, si è infilata un tappo così a fondo nell'orecchio che adesso ha male al timpano.

Scendiamo per la colazione e l'accoglienza è calorosa come sempre: ci portano subito caffè, succo e un piatto di frutta fresca. Ordiniamo due piatti di Banana Crab aspettandoci i pancake del pranzo di ieri, ma invece ci arrivano due crepes di banane. Poco male, le divoriamo lo stesso.

Mentre saldiamo il conto con Oshi arriva anche Ronnie a chiederci come stiamo e a informarsi se abbiamo un passaggio disponibile per andare via, subito pronto a offrire aiuto in caso di necessità. In tutto questo ci chiedono almeno tre volte se vogliamo altro caffè o tè. Questa frase ce l'avranno ripetuta una ventina di volte in questi giorni, e fa sorridere se si pensa che caffè e tè sono gratis — hanno persino una sezione apposita "FREE" sul menù.

Alle 9:30 puntuali li salutiamo per avviarci verso l'ingresso del villaggio, ma ovviamente Oshi si fa avanti per accompagnarci. È così gentile e premuroso che, per prenderlo in giro un po', gli chiedo se posso portargli un caffè.

Non avevamo notizie dell'autista e non sapevamo se si sarebbe presentato puntuale o se avrebbe trovato facilmente il posto, ma mentre attraversiamo il villaggio eccolo che ci viene incontro. Carichiamo la macchina, salutiamo e partiamo.

Poco lontano dal villaggio l'autista si ferma e una donna con una bambina per mano gli consegna un borsello: ci dice che sono moglie e figlia e che aveva dimenticato il portafoglio. Per questo trasferimento mi ero messo d'accordo con il resort dell'isola di Gili Asahan, la nostra destinazione del giorno, e mi sorprende che l'autista viva proprio lì vicino — ma questo spiega anche come abbia trovato il villaggio così facilmente.

Il nostro autista è simpatico ma non ha un ottimo inglese, e noi siamo stanchi e curiosi di quanto vediamo lungo i bordi della strada, quindi non parliamo tantissimo. Gli faccio però qualche domanda ed è così che mi racconta della differenza tra il tabacco locale e quello per Philip Morris, e mi spiega perché vediamo ovunque depositi di gusci di cocco e mattoni: i gusci vengono fatti bruciare e usati per cuocere i mattoni.

Lungo la strada vediamo tante scuole con i cortili pieni di scolaresche che marciano all'unisono. Sono principalmente bambine e vorrei chiedere all'autista il motivo, ma perdo l'occasione: dovrò scoprirlo in altro modo. Elisa scatta qualche foto dal finestrino e fa qualche video, poi come suo solito si addormenta e recupera un po' di sonno. Peggio per lei stavolta, perché mentre saliamo per dei tornanti incrociamo una famiglia di una decina di scimmie grigie — con i piccoli appesi al ventre delle mamme — che si spostano in fila indiana lungo il guardrail. Purtroppo non riesco a scattare la foto, così rimarrà solo nella mia memoria.

In poco meno di tre ore, con una sosta bagno nella peggior latrina d'Indonesia su cui preferisco non soffermarmi, giungiamo a Pantai Kores, da dove ci aspetta il traghetto per Gili Asahan. Secondo il mio itinerario il traghetto sarebbe dovuto partire tra un'ora, ma le cose vanno decisamente meglio delle aspettative per due motivi: uno, appena arriviamo ci imbarcano immediatamente insieme a una sola famiglia di olandesi; due, non si tratta di un traghetto ma di un catamarano monomotore da pesca. È efficientissimo e in un quarto d'ora siamo sull'isola.

Il Pearl Beach Resort batte ogni aspettativa, così come il posto. Ovunque sulla spiaggia bianca ci sono frammenti di conchiglie e coralli, il mare è trasparente e una brezza leggera rende il sole piacevole. Dietro la spiaggia, tra le palme di cocco, una serie di lettini con materasso sono disposti a coppie ben distanziate per lasciare agli ospiti la privacy. C'è un palco rialzato per rilassarsi o scrutare l'orizzonte, una zona yoga con materassini, un'area relax con divanetti, tavolini e una mini biblioteca multilingue. La reception si trova tra il ristorante all'aperto e il centro di scuba diving, dove è possibile organizzare immersioni e noleggiare attrezzature e kayak.

In attesa del check-in siamo accolti al ristorante con un drink di benvenuto a base di acqua, limone e canna da zucchero. Il proprietario si presenta — notiamo che ha un'unghia lunghissima — e ci chiede da dove veniamo e dove siamo stati finora. Quando gli diciamo di Al Sasaki ci sorprende: conosce bene Ronnie, ed è stato proprio lui a insegnargli a fare turismo. Si conoscono dal 1999, da quando una camera da Ronnie costava 7.500 rupie indonesiane, ossia cinquanta centesimi a notte.

La nostra stanza non è ancora pronta così ne approfittiamo per pranzare: Young Papaya Chicken Salad per me e un Nasi Goreng with Satay Chicken per Elisa, riso fritto avvolto in una frittata con verdure e pollo satay. Io non ci ho proprio azzeccato, ma almeno mi sono tenuto leggero.

Dopo pranzo la stanza è pronta e posso descrivere il resto del resort. Sul retro si trovano tanti lodge indipendenti che costituiscono le camere degli ospiti. Ognuno ha un piccolo patio con un paio di sedie, un tavolino e un'enorme amaca di tela, su cui bisogna lasciare scarpe e ciabatte prima di entrare. Le porte sono a vetri scorrevoli, il tetto è di paglia; all'interno si trova un letto a baldacchino e delle mensole a forma di prua di barca, tutto interamente in legno. Il bagno è sul retro, all'aperto, con doccia in pietra e porta-sapone di bambù su vassoi di madreperla. Ogni oggetto della stanza e del patio è decorato con fiori di Tiarè.

Come benvenuto ci hanno lasciato una caraffa d'acqua con frutta: banane indonesiane, mandarini e salak, un frutto locale con la buccia a scaglie come un serpente, dentro ha spicchi che somigliano a quelli dell'aglio e un sapore dolce e fibroso. La camera è favolosa. Ricorda quella del Messico, la capanna nella giungla, ma probabilmente questa vince. E poi costa un quarto, quindi sì, stravince.

Siamo euforici, abbiamo ancora il pomeriggio davanti e vorremmo esplorare, nuotare e fare mille cose, ma siamo anche molto stanchi quindi decidiamo di rilassarci per il resto della giornata. Al ristorante ci consegnano gli asciugamani e andiamo a cercare un posto dove coricarci: io mi sistemo sotto una tenda di foglie di palma, Eli si corica su un'amaca. Resiste forse cinque minuti — senza cadere — poi va sulla spiaggia dove c'è più sole.

Io provo a dormire: non ho mai dormito in spiaggia, non ci sono mai riuscito, ma questa potrebbe essere l'eccezione visto che sono stanco morto e il materasso è davvero comodo. Anche stavolta però non ci sono sorprese e dopo un po' rinuncio e mi metto a leggere.

Quando si fanno le quattro mi ricordo che al check-in ci avevano spiegato che tra le quattro e le cinque vengono serviti tè e caffè gratis al ristorante, così decido di andare a servirmi. Provo a chiamare Elisa ma dorme beata sulla spiaggia: bocca aperta, respiro pesante, manca solo il rivolino di saliva a completare la scena. Non la disturbo e vado da solo.

Mi servo un caffè bollente e prendo un dolcetto — la prima cosa non particolarmente buona che assaggio qui, se si esclude la cena la prima sera a Jakarta — e mi sistemo a un tavolino d'angolo a leggere con il mare accanto. Qualcuno conversa con i locali al bancone, qualcuno fa amicizia. Io ho il mio libro, che è un linguaggio universale per comunicare "non rompetemi i coglioni", e in effetti funziona. Anche il caffè funziona: mi toglie la sonnolenza e fa sparire quel lieve mal di testa.

Terminato il break vado a trovare Elisa, che nel frattempo si è svegliata e si aggira per la spiaggia alla ricerca di conchiglie: mani dietro la schiena, testa china verso il basso, piccoli passi — mi ricorda un anziano davanti a un cantiere. Le ricordo che ci hanno detto di non portare via nulla né dal mare né dalla spiaggia, non tanto per salvaguardare i coralli locali quanto per non riempirmi casa di conchiglie. Funziona: butta via quelle che aveva preso, ma sono sicuro che non la scamperò per tutta la nostra permanenza sull'isola.

Passeggiamo insieme lungo la spiaggia e raggiungiamo una zona di pescatori da cui godiamo di un tramonto dietro le montagne della vicina Lombok. Il sole è rosso fuoco.

Torniamo in camera a prepararci mentre inizia a fare buio e ci docciamo per andare a cena. Al ristorante hanno messo una musichetta jazz di sottofondo, perfetta per l'atmosfera. Ci accomodiamo insieme ai pochi altri ospiti a piedi scalzi al tavolo più vicino al mare, dove sentiamo le onde infrangersi sulla spiaggia a pochi metri da noi.

Decido di rifarmi dopo il pranzo mediocre e ordino un Mie Goreng vegetariano, noodles con verdure, mentre Elisa ordina un Chicken Pelalah Curry, filetto di pollo in tempura di mais con salsa locale, croccante fuori e morbido dentro. Da bere per me una pilsner indonesiana, la Bintang, e per Elisa acqua di cocco direttamente dal cocco, con una cannuccia di acciaio. Niente plastica qui.

Tutto gustosissimo, ma abbiamo ancora fame — soprattutto Elisa — così lei ordina patate dolci con salsa come dessert e io un waffle fatto in casa. Le sue patatine sono buonissime, prodotti locali davvero genuini e sani.

Terminata la cena andiamo a coricarci sulle sdraio e ammiriamo le costellazioni sopra di noi: senza illuminazione artificiale qui si vedono benissimo e riusciamo addirittura a scattare qualche foto delle stelle. La giornata ce la siamo portata a casa anche stavolta e possiamo tornare in camera a dormire, ma non prima di aver verificato due cose: nei paraggi non ci sono né galli né moschee.


\pagestyle{empty} % disable numbers

%coppia
\portraitFullscreen{images/day4/PXL_20240812_002552051.jpg}
\threeImagesLargeMargins{images/day4/PXL_20240812_054618071.jpg}
                        {images/day4/PXL_20240812_054552127.jpg} 
                        {images/day4/PXL_20240812_054602592.jpg}
%

%coppia
\twoImagesLargeMargins{images/day4/IMG_3950.JPG}
                      {images/day4/IMG_4000.JPG}
\twoImagesLargeMargins{images/day4/IMG_3988.JPG}
                      {images/day4/IMG_3982.JPG}
%

%coppia
\splitImageLeft{images/day4/IMG_3962.JPG}
\splitImageRight{images/day4/IMG_3962.JPG}
%

%coppia
\splitImageLeft{images/day4/PXL_20240812_051118946.jpg}
\splitImageRight{images/day4/PXL_20240812_051118946.jpg}
%

%coppia
\landscapeThinMarginCentered{images/day4/PXL_20240812_061522252.jpg}
\landscapeThinMargin{images/day4/6F2D890C-2BED-409E-B401-CB530041C1DB_1_201_a.jpeg}
%

%coppia
\splitImageLeft{images/day4/PXL_20240812_072152483.jpg}
\splitImageRight{images/day4/PXL_20240812_072152483.jpg}
%

%coppia
\threeImagesThinMargins{images/day4/PXL_20240812_070437859.jpg}
                        {images/day4/PXL_20240812_070514381.MP.jpg}
                        {images/day4/PXL_20240812_070522803.MP.jpg}
\twoImagesLargeMargins{images/day4/PXL_20240812_072056584.jpg}
                      {images/day4/PXL_20240812_072018475.MP.jpg}

%
\twoImagesLargeMargins{images/day4/PXL_20240812_085426388.MP.jpg}
                      {images/day4/PXL_20240812_091634718.jpg}
\twoImagesLargeMargins{images/day4/PXL_20240812_101609106.MP.jpg}
                      {images/day4/PXL_20240812_100145461.MP.jpg}
%

%coppia
\twoImagesLargeMargins{images/day4/PXL_20240812_095506076.jpg}
                      {images/day4/PXL_20240812_095636039.MP.jpg}
\twoImagesLargeMargins{images/day4/PXL_20240812_100421105.MP.jpg}
                      {images/day4/PXL_20240812_100349730.MP.jpg}
%

%coppia
\splitImageLeft{images/day4/PXL_20240812_100250360.MP.jpg}
\splitImageRight{images/day4/PXL_20240812_100250360.MP.jpg}
%

%coppia
\landscapeLargeMargin{images/day4/PXL_20240812_120842708.jpg}
\portraitFullscreen{images/day4/PXL_20240812_133913430.NIGHT.jpg}
%

\portraitFullscreen{images/day4/IMG_3963.JPG}

\pagestyle{fancy} % enable numbers
